\documentclass[main.tex]{subfiles}
\begin{document}

\section{Exercise 6: CMRR Improvement in Differential Amplifier}

\begin{figure}[H]
    \centering
    \includegraphics[width=0.5\linewidth]{images/exercise6_circuit.png}
    \caption{CMOS differential amplifier}
\end{figure}

CMRR = $|A_d/A_{cm}|$ tells you how well the circuit amplifies differential signals versus common-mode stuff. For single-ended output, $A_{cm} \approx -R_D/(2R_B)$ which gives CMRR $\approx g_m R_B$ - so it's basically proportional to the tail impedance.

\textbf{(a) CMRR improvement techniques:}

\textit{Cascode tail current source:} If you cascode MB, the tail resistance goes up to $g_m r_{ds}^2$, which can give you 20-40 dB better CMRR just from the higher tail impedance.

\textit{Larger M1-M2 dimensions:} Making W/L bigger improves matching since $\sigma(\Delta V_T) \propto 1/\sqrt{WL}$. This reduces how much common-mode signal gets converted to differential because of mismatch.

\textit{Cascode load:} You can also cascode M3-M4 which increases both the differential gain $A_d$ and the output impedance, helping CMRR indirectly.

\textit{Bypass capacitor:} Adding a cap to AC ground the tail node helps reject AC common-mode signals.

\textbf{(b) Trade-offs:}

The cascode tail gives you a lot of CMRR improvement and is pretty straightforward to implement, but you need $2V_{DS,sat}$ headroom which cuts into your input common-mode range. This can be a problem in low-voltage designs.

Making the input transistors bigger gives better matching and noise, but you use more area and get higher parasitics which might hurt your bandwidth.

Cascode load gets you higher gain and CMRR but reduces the output swing. You also need extra bias generators and the high-impedance nodes can cause stability issues.

The bypass cap is simple and cheap for rejecting AC common-mode, but you need a big capacitor if you want it to work at low frequencies. Also it does nothing for DC CMRR.

\end{document}
